\documentclass[a4paper,11pt]{article}

\usepackage{fullpage}
\usepackage[utf8]{inputenc}
\usepackage[T1]{fontenc}
\usepackage{amsmath,amsthm}
\usepackage{amsfonts}
\usepackage{graphicx}
\usepackage{latexsym}
\usepackage{float}
\usepackage{comment}
\usepackage{array}
\usepackage{booktabs}

% --- Packages added for pedagogical enrichment ---
\usepackage{xcolor}
\usepackage{listings}
\usepackage[most]{tcolorbox}
\usepackage{hyperref}

%%%%%%%%%%%%%%% LISTINGS CONFIGURATION %%%%%%%%%%%%%%%%%%%

\definecolor{codegreen}{rgb}{0.25,0.49,0.25}
\definecolor{codegray}{rgb}{0.5,0.5,0.5}
\definecolor{codepurple}{rgb}{0.58,0,0.82}
\definecolor{codered}{rgb}{0.7,0.1,0.1}
\definecolor{codeblue}{rgb}{0.0,0.0,0.7}
\definecolor{backcolour}{rgb}{0.97,0.97,0.97}

\lstset{
    language=PHP,
    backgroundcolor=\color{backcolour},
    commentstyle=\color{codegreen}\itshape,
    keywordstyle=\color{codeblue}\bfseries,
    stringstyle=\color{codered},
    basicstyle=\ttfamily\small,
    breakatwhitespace=false,
    breaklines=true,
    captionpos=b,
    keepspaces=true,
    numbers=left,
    numbersep=8pt,
    numberstyle=\tiny\color{codegray},
    showspaces=false,
    showstringspaces=false,
    showtabs=false,
    tabsize=4,
    frame=single,
    rulecolor=\color{codegray},
    xleftmargin=1.5em,
    framexleftmargin=1.5em,
    literate={é}{{\'e}}1 {è}{{\`e}}1 {ê}{{\^e}}1 {à}{{\`a}}1 {ù}{{\`u}}1 {î}{{\^i}}1 {ô}{{\^o}}1
}

%%%%%%%%%%%%%%% TCOLORBOX ENVIRONMENTS %%%%%%%%%%%%%%%%%%

% Course reminder (light blue background)
\newtcolorbox{rappel}[1][]{
    colback=blue!5!white,
    colframe=blue!50!black,
    fonttitle=\bfseries,
    title={\large Course Reminder},
    #1
}

% Key takeaway (light green background)
\newtcolorbox{pointcle}[1][]{
    colback=green!5!white,
    colframe=green!40!black,
    fonttitle=\bfseries,
    title={Key Takeaway},
    #1
}

% Warning (light orange background)
\newtcolorbox{attention}[1][]{
    colback=orange!8!white,
    colframe=orange!60!black,
    fonttitle=\bfseries,
    title={Warning},
    #1
}

% Ethical disclaimer (red background)
\newtcolorbox{ethique}[1][]{
    colback=red!5!white,
    colframe=red!60!black,
    fonttitle=\bfseries\large,
    title={Ethical and Legal Disclaimer},
    #1
}

%%%%%%%%%%%%%%% COMMANDS %%%%%%%%%%%%%%%%%%%%%%%%%%%

\newcommand{\itemb}{\item[$\bullet$]}

\newcommand{\Fd}{\mathbb{F}}
\newcommand{\Znat}{\mathbb{Z}}
\newcommand{\Nnat}{\mathbb{N}}

%%%%%%%%%%%%%%% ENVIRONMENTS %%%%%%%%%%%%%%%%%%%%%%%

\theoremstyle{plain}
\newtheorem{lemme}{Lemma}
\newtheorem{algorithme}{Algorithm}
\newtheorem{corolaire}{Corollary}
\newtheorem{propriete}{Property}
\newtheorem{proposition}{Proposition}
\newtheorem{theoreme}{Theorem}
\newtheorem{definition}{Definition}

\theoremstyle{definition}
\newtheorem{exercice}{Exercise}
\newtheorem{remarque}{Remark}
\newtheorem{exemple}{Example}

\hypersetup{
    colorlinks=true,
    linkcolor=blue!60!black,
    urlcolor=blue!60!black
}

%%%%%%%%%%%%%% DOCUMENT %%%%%%%%%%%%%%%%%%%%%%%%%%%%


\begin{document}

\hbox{\raisebox{0.4em}{\vrule depth 0pt height 0.5pt width \textwidth}}
\noindent \textbf{University of Perpignan} \hfill  L3 Computer Science \\
 \hfill C. Legrand \hfill  Year \the\year \\
\smallskip
\begin{center}
{
  {\scshape \large Lab Work 7 -- Security}  \\
 Code Injection
}
\end{center}
\hbox{\raisebox{0.4em}{\vrule depth 0pt height 0.9pt width \textwidth}}

\bigskip

% --- Ethical disclaimer ---
\begin{ethique}
The techniques presented in this lab work are studied within a \textbf{strictly academic context} in order to understand the mechanisms used by attackers and thus be better able to defend against them.

\textbf{Any use of these techniques outside the educational context is illegal} and punishable by criminal penalties (Chapter III of the French Penal Code, Articles 323-1 to 323-8 -- \textit{Offences against Automated Data Processing Systems}):
\begin{itemize}
    \item Unauthorized access to or remaining in an ADPS (Art.~323-1): \textbf{3 years} imprisonment and \textbf{\texteuro{}100,000} fine
    \item Obstructing the operation of an ADPS (Art.~323-2): \textbf{5 years} imprisonment and \textbf{\texteuro{}150,000} fine
    \item Fraudulent introduction of data into an ADPS (Art.~323-3): \textbf{5 years} imprisonment and \textbf{\texteuro{}150,000} fine
\end{itemize}
Penalties increased to \textbf{7 years} and \textbf{\texteuro{}300,000} when the targeted system processes personal data operated by the State, and up to \textbf{10 years} and \textbf{\texteuro{}300,000} for organized crime (Art.~323-4-1). \textit{Penalties updated by the LOPMI law no.~2023-22 of January 24, 2023.}
\end{ethique}

\bigskip

\begin{tcolorbox}[colback=gray!5!white, colframe=gray!60!black, title={Setup with Docker}]
For this lab work, the environment (Apache + PHP + MySQL server) is provided as Docker containers. No virtual machine is needed.

\medskip
\textbf{Prerequisite: install Docker}

\medskip
\textit{On Ubuntu/Debian:}
\begin{lstlisting}[language=bash,numbers=none]
# Install Docker
sudo apt update
sudo apt install -y docker.io docker-compose-v2
# Add your user to the docker group
sudo usermod -aG docker $USER
# Log out and log back in for the group
# change to take effect
\end{lstlisting}

\textit{On macOS/Windows:} install \href{https://www.docker.com/products/docker-desktop/}{Docker Desktop}.

\medskip
\textbf{Starting the environment}

\begin{lstlisting}[language=bash,numbers=none]
# Navigate to the lab directory
cd travaux-diriges/04-injection-code/td07/

# Start the containers (first time: ~2 min)
docker compose up -d

# Check that everything is running
docker compose ps
\end{lstlisting}

On first launch, Docker will:
\begin{itemize}
    \item build an Apache+PHP container with the \texttt{mysqli} extension,
    \item start a MySQL 8.0 container with the \texttt{l2info} database,
    \item automatically load the tables and data from \texttt{bd.sql}.
\end{itemize}

\medskip
\textbf{Accessing the web pages}

The lab files are accessible in your browser at:

\begin{center}
\texttt{http://localhost:8080/td7/}
\end{center}

For example: \texttt{http://localhost:8080/td7/page-5-commentaires.php}

\medskip
\textbf{Stopping and cleaning up}

\begin{lstlisting}[language=bash,numbers=none]
# Stop the containers (data is preserved)
docker compose stop

# Restart
docker compose start

# Stop and remove everything (containers + data)
docker compose down -v
\end{lstlisting}

\medskip
\begin{attention}
In the PHP files, the MySQL server is accessible via the hostname \texttt{"db"} (the Docker service name), not \texttt{"localhost"}. This is normal Docker behavior: each service runs in a separate container, identified by its name.
\end{attention}

\end{tcolorbox}

\bigskip


%%%%%%%%%%%%%%%%%%%%%%%%%%%%%%%%%%%%%%%%%%%%%%%%%%%%%%%%%%%%%%%%%%
%% EXERCISE 1: HTML Form Manipulation
%%%%%%%%%%%%%%%%%%%%%%%%%%%%%%%%%%%%%%%%%%%%%%%%%%%%%%%%%%%%%%%%%%

\begin{exercice}[HTML Form Manipulation]

Consider a web shopping application whose HTML form is as follows (file \texttt{achat-iphone5.html}):

\begin{lstlisting}[language=HTML]
<html>
<body>

<form method="post" action="enregistre-achat.php">
Produit: iphone5 <br/>
Prix : 500 euros <br/>
Quantit&eacute <input type="text" name ="quantite">
  (Quantite max est 90) <br/>
<input type="hidden" name="prix" value="500" >
<input type="submit" value="buy">
</form>

</body>
</html>
\end{lstlisting}

The PHP script \texttt{enregistre-achat.php} records the purchase in the database using the values sent by the form.

\begin{enumerate}
    \item Analyze the HTML form. Identify the vulnerability related to how the price is sent to the server.
    \item Propose a method to buy an iPhone for \texteuro{}5 instead of \texteuro{}500.
    \item Propose a server-side fix to prevent this attack.
\end{enumerate}

\begin{pointcle}
Hidden fields and any client-side data can be modified by the user. Validation must always be performed server-side.
\end{pointcle}

\end{exercice}


%%%%%%%%%%%%%%%%%%%%%%%%%%%%%%%%%%%%%%%%%%%%%%%%%%%%%%%%%%%%%%%%%%
%% EXERCISE 2: Cross-Site Scripting (XSS)
%%%%%%%%%%%%%%%%%%%%%%%%%%%%%%%%%%%%%%%%%%%%%%%%%%%%%%%%%%%%%%%%%%

\begin{exercice}[Cross-Site Scripting (XSS)]

\begin{rappel}
There are three types of XSS:
\begin{itemize}
    \item \textbf{Reflected XSS}: the payload is included in the request (URL, GET/POST parameter) and immediately reflected back in the response without being stored.
    \item \textbf{Stored XSS}: the payload is saved on the server (database, file) and displayed to all subsequent visitors of the page.
    \item \textbf{DOM-based XSS}: the payload is interpreted directly by client-side JavaScript, without passing through the server.
\end{itemize}
\end{rappel}

\medskip

\textbf{a) Stored XSS -- JavaScript alert}

\medskip

Consider the page \texttt{page-5-commentaires.php} whose code is as follows:

\begin{lstlisting}[language=PHP]
<html>
<H1> Computer Security  </H1>
<IMG src="img_security.jpeg"> Computer security...
<BR/><hr/>
<h3> Visitors comments :</h3>
<?php
$conn = new mysqli("db", "l2info", "l2info", "l2info");
if ($conn->connect_error) {
  die("Connection failed: " . $conn->connect_error);
}

if( isset($_POST['login']) && isset($_POST['commentaire']) )
{
    $sql = 'INSERT INTO Commentaires(login,commentaire)
      VALUES("' . $_POST['login']. '","'
      . $_POST['commentaire'].'");';
    $result0 = $conn->query($sql);
}

$sql = 'SELECT id,login, commentaire FROM Commentaires
  ORDER BY id DESC LIMIT 4 ;';
$result = $conn->query($sql);

if ($result->num_rows > 0) {
    echo "<ul>";
    while($row = $result->fetch_assoc()) {
        echo "<li>";
        echo "[" . $row["login"]. "]: "
          . $row["commentaire"]."<br>";
        echo "</li>";
    }
    echo "</ul>";
} else {
    echo "No comments.";
}
$conn->close();
?>
<BR/><hr/><BR/>
<form name="postcomment" method="POST"
  action="page-5-commentaires.php">
Pseudo: <input name="login"> <BR/>
Comment: <input type="text" name="commentaire"><BR/>
<input type="submit">
</form>
</html>
\end{lstlisting}

\begin{enumerate}
    \item Explain what this page does.
    \item Propose an injection in the comment field so that a JavaScript alert is displayed when the page loads.
\end{enumerate}

\medskip

\textbf{b) Stored XSS -- defacement}

\medskip

Using the same page, inject fraudulent HTML to make a fake ``References'' section appear with a link to a fictitious website.

\medskip

\textbf{c) Reflected XSS via GET}

\medskip

Consider the following PHP code:

\begin{lstlisting}[language=PHP]
<html>
<body>
<?php
if( isset($_GET['UserInput'])){
  $out ='you entered: "'
    . $_GET['UserInput'] . '".';
}
else {
  $out = '<form method="GET">
    enter something here:';
  $out .= '<input name="UserInput" size="50">';
  $out .= '<input type="submit">';
  $out .= '</form>';
}
print $out;
?>
</body>
</html>
\end{lstlisting}

Submit an input that triggers a JavaScript alert in the result.

\medskip

\textbf{d) Reflected XSS via POST}

\medskip

Now consider the following code:

\begin{lstlisting}[language=PHP]
<html>
<body>
<?php
if( isset($_POST['UserInput'])){
  $out ='you entered: "'
    . $_POST['UserInput'] . '".';
}
else {
  $out='No data was sent to the script';
}
print $out;
?>
</body>
</html>
\end{lstlisting}

Build an external HTML page containing a hidden POST form that automatically submits an XSS payload to this page.

\medskip

\textbf{e) DOM-based XSS}

\medskip

Consider the following HTML page:

\begin{lstlisting}[language=HTML]
<html>
<body>
<div id="welcome"></div>
<script>
  var name = document.location.hash.substring(1);
  document.getElementById("welcome").innerHTML
    = "Welcome " + name;
</script>
</body>
</html>
\end{lstlisting}

\begin{enumerate}
    \item Craft a URL that exploits this vulnerability to execute JavaScript.
    \item Explain why this type of XSS is harder to detect server-side.
\end{enumerate}

\begin{attention}
Stored XSS is the most dangerous because it affects all visitors to the page, without requiring a crafted link.
\end{attention}

\end{exercice}


%%%%%%%%%%%%%%%%%%%%%%%%%%%%%%%%%%%%%%%%%%%%%%%%%%%%%%%%%%%%%%%%%%
%% EXERCISE 3: SQL Injection
%%%%%%%%%%%%%%%%%%%%%%%%%%%%%%%%%%%%%%%%%%%%%%%%%%%%%%%%%%%%%%%%%%

\begin{exercice}[SQL Injection]

\begin{rappel}
\textbf{Difference between XSS and SQL injection:}
\begin{itemize}
    \item \textbf{XSS} injects malicious code into the \textit{victim's browser} (client-side). The attack targets the website's users.
    \item \textbf{SQL injection} injects malicious code into the \textit{server's database} (server-side). The attack directly targets the site's data.
\end{itemize}
\end{rappel}

\medskip

\textbf{a) Analyzing the vulnerable code}

\medskip

Consider the page \texttt{page-transaction-sql.php}:

\begin{lstlisting}[language=PHP]
<html>
<?php
$servername = "db";
$username = "l2info";
$password = "l2info";
$dbname = "l2info";

$conn = new mysqli($servername, $username,
  $password, $dbname);
if ($conn->connect_error) {
  die("Connection failed: "
    . $conn->connect_error);
}
echo 'Name: "' . $_POST['UserInput'].'"<BR/>';
$sql = 'SELECT montant FROM Transaction WHERE nom="'
  . $_POST['UserInput']  . '";';
echo "SQL Query=" . $sql . "<BR/>";
$result = $conn->query($sql);
echo "Result:<BR/>";
if ($result->num_rows > 0) {
  while($row = $result->fetch_assoc()) {
    echo "Name: ". $_POST['UserInput']
      ." Amount:" . $row["montant"]."<br/>";
  }
} else {
  echo "0 results";
}
$conn->close();
?>
</html>
\end{lstlisting}

And the form \texttt{injection-sql.html}:

\begin{lstlisting}[language=HTML]
<html>
<body>
<form name="exploit-sql" method="POST"
  action="http://localhost:8080/td7/page-transaction-sql.php">
<input name="UserInput">
<input type="submit">
</form>
</body>
</html>
\end{lstlisting}

\begin{enumerate}
    \item Explain how the SQL query is built. Identify the vulnerability.
\end{enumerate}

\medskip

\textbf{b) Display the entire table}

\medskip

Propose an injection to display the entire \texttt{Transaction} table. Explain the resulting query step by step.

\medskip

\textbf{c) Attempting a DROP TABLE}

\medskip

Attempt to delete the \texttt{Test} table via SQL injection. Explain why this may fail with MySQL.

\medskip

\textbf{d) UNION-based}

\medskip

Extract the data from the \texttt{Test} table (column \texttt{val}) using a \texttt{UNION SELECT} injection.

\textit{Hint:} start by determining the number of columns in the original \texttt{SELECT} by using \texttt{ORDER BY 1}, \texttt{ORDER BY 2}, etc. until you get an error. Then construct the \texttt{UNION SELECT} with the correct number of columns.

\medskip

\textbf{e) Blind SQL injection}

\medskip

Now suppose the page no longer displays the query results but only shows ``Transaction found'' or ``No transaction''.

\begin{enumerate}
    \item How could you determine whether the MySQL version is 5.x? (concept of \textit{boolean-based blind SQL injection})
    \item What would be the principle for extracting data character by character?
\end{enumerate}

\begin{pointcle}
SQL injection works because the query mixes SQL code and user data in the same string. The solution is to separate them using \textbf{prepared statements}.
\end{pointcle}

\end{exercice}


%%%%%%%%%%%%%%%%%%%%%%%%%%%%%%%%%%%%%%%%%%%%%%%%%%%%%%%%%%%%%%%%%%
%% EXERCISE 4: Countermeasures
%%%%%%%%%%%%%%%%%%%%%%%%%%%%%%%%%%%%%%%%%%%%%%%%%%%%%%%%%%%%%%%%%%

\begin{exercice}[Countermeasures]

\textbf{a) Fixing XSS}

\medskip

Modify the code of \texttt{page-5-commentaires.php} to prevent XSS injection. Which PHP function escapes special HTML characters? Provide the corrected code for the comment display line.

\medskip

\textbf{b) Fixing SQL injection}

\medskip

Rewrite \texttt{page-transaction-sql.php} using PDO prepared statements instead of \texttt{mysqli} concatenation. Provide the complete corrected code.

\medskip

\textbf{c) Additional defenses}

\begin{enumerate}
    \item What is the \texttt{Content-Security-Policy} HTTP header and how does it protect against XSS?
    \item What is an \texttt{HttpOnly} cookie and why is it useful?
    \item Why is the principle of least privilege important for the database?
\end{enumerate}

\begin{pointcle}
Defense in depth combines multiple mechanisms: output escaping, prepared statements, CSP, \texttt{HttpOnly} cookies, input validation.
\end{pointcle}

\end{exercice}


\end{document}
